% Discussion paragraph (≈250w) — epistemic status of REAL-flagged motifs.
% For task t-mozv8oul. Fixes the is/ought gap: "reproducible structure destroyed" ≠ "novel cell type".
% Guards the team from over-claiming. Anchors: motif $w \in \mathcal{W}$; LRS(w); RareShield; Hacking-style intervention realism.

\paragraph*{What a REAL flag is, and is not.}
A motif flagged by REAL is the claim that a local, reproducible geometric structure---witnessed across independent batches before integration, with measurable local-mass distortion after---has been erased by the integration map beyond what admissible batch-effect deformations can explain. This is a claim about \emph{evaluation}, not about ontology. It is not a claim that a new cell type exists. A motif $w \in \mathcal{W}$ with low $\mathrm{LRS}(w)$ licenses exactly one inference: the integration has discarded reproducible structure of the kind a biologist would want to interrogate. Whether that structure corresponds to a novel cell type, a transitional state along an already-known trajectory, a rare disease-specific configuration, or a technical sub-mode that merely looks biological, is a further question that REAL does not answer. Upgrading a motif to a cell-type claim requires the usual burden of single-cell biology: independent sampling, orthogonal markers, targeted protocol (sorting, spatial validation, experimental perturbation), and community replication. REAL's value is that it identifies \emph{where to look}, not what has been found. Conflating the two would re-create the very error this paper is designed to expose: treating a framework-internal signal as a fact about the world. Our strong claim is therefore deliberately asymmetric: a high motif set $|\mathcal{W}|$ under a candidate integration is evidence of measurement insufficiency; a low motif set is evidence that the integration respected the syntactic reproducibility criterion. Ontology is downstream of both.
